\chapter{Программный код\label{appendix-extra-examples}}							% 

\begin{lstlisting}[language=Java, breaklines=true]
package com.tcp.algorithm;

import com.tcp.model.TestCase;
import java.util.*;


/**
 * базовый жадный алгоритм
 * на каждом шаге выбирает тест, покрывающий наибольшее количество
 */

public class AdditionalPrioritizer implements Prioritizer {

    @Override
    public List<TestCase> prioritize(List<TestCase> testPool) {
        List<TestCase> result = new ArrayList<>();
        Set<String> coveredMethods = new HashSet<>();
        //копия пула, чтобы не модифицировать исходный список
        List<TestCase> remaining = new ArrayList<>(testPool);

        while (!remaining.isEmpty()) {
            TestCase bestTest = null;
            int maxNewCoverage = -1;

            //поиск теста с максимальным новым покрытием
            for (TestCase test : remaining) {
                int newCount = 0;
                for (String method : test.getCoveredMethods()) {
                    if (!coveredMethods.contains(method)) {
                        newCount++;
                    }
                }

                if (newCount > maxNewCoverage) {
                    maxNewCoverage = newCount;
                    bestTest = test;
                } else if (newCount == maxNewCoverage && bestTest != null) {
                    //если покрытие одинаковое берем тест с большей fault proneness
                    if (test.getFaultProneness() > bestTest.getFaultProneness()) {
                        bestTest = test;
                    }
                }
            }
            //fallback: если все оставшиеся тесты не покрывают ничего нового
            if (bestTest == null) {
                bestTest = remaining.get(0);
            }

            result.add(bestTest);
            remaining.remove(bestTest);
            coveredMethods.addAll(bestTest.getCoveredMethods());

        }
        return result;
    }

    @Override
    public String getName() {
        return "Additional";
    }

}
\end{lstlisting}


\begin{lstlisting}[language=Java, breaklines=true]
package com.tcp.algorithm;

import com.tcp.fitness.GaFitnessFunction;
import com.tcp.metric.ApfdCalculator;
import com.tcp.model.TestCase;

import java.util.*;

/**
 * Генетический алгоритм приоритизации тестов.
 *
 * Поддерживает два режима фитнесс-функции:
 *   1. Прямой APFD -- если переданы testToFaults и totalFaults (как в статье Paygude et al.)
 *   2. GaFitnessFunction -- взвешенная сумма coverage + faultProneness + diversity
 *
 * Операторы:
 *   инициализация: случайное перемешивание
 *   селекция: турнирная (tournament selection)
 *   кроссовер: OX (Order Crossover)
 *   мутация: адаптивное число swap = sqrt(n)
 *   элитизм: лучшая особь переходит без изменений
 */
public class GeneticPrioritizer implements Prioritizer {

    private static final int TOURNAMENT_SIZE = 3;

    private final int populationSize;
    private final int generations;
    private final double crossoverRate;
    private final double mutationRate;
    private final GaFitnessFunction fitnessFunction; // null если используется APFD-режим
    private final Map<String, Set<String>> testToFaults; // null если используется GaFitnessFunction
    private final int totalFaults;
    private final Random rng = new Random(42);

    // Конструктор для режима GaFitnessFunction (coverage + faultProneness + diversity)
    public GeneticPrioritizer(int populationSize, int generations,
                              double crossoverRate, double mutationRate,
                              GaFitnessFunction fitnessFunction) {
        this(populationSize, generations, crossoverRate, mutationRate, fitnessFunction, null, 0);
    }

    // Конструктор для режима прямого APFD (как в статье Paygude et al.)
    public GeneticPrioritizer(int populationSize, int generations,
                              double crossoverRate, double mutationRate,
                              GaFitnessFunction fitnessFunction,
                              Map<String, Set<String>> testToFaults,
                              int totalFaults) {
        this.populationSize = populationSize;
        this.generations = generations;
        this.crossoverRate = crossoverRate;
        this.mutationRate = mutationRate;
        this.fitnessFunction = fitnessFunction;
        this.testToFaults = testToFaults;
        this.totalFaults = totalFaults;
    }

    // Выбирает режим фитнесса: прямой APFD или GaFitnessFunction
    private double fitness(List<TestCase> seq) {
        if (testToFaults != null && totalFaults > 0)
            return ApfdCalculator.calculate(seq, testToFaults, totalFaults);
        return fitnessFunction.calculateFitness(seq);
    }

    @Override
    public List<TestCase> prioritize(List<TestCase> testPool) {
        if (testPool.isEmpty()) return new ArrayList<>();

        // адаптивная интенсивность мутации: sqrt(n) swap
        final int mutationIntensity = Math.max(1, (int) Math.sqrt(testPool.size()));

        List<List<TestCase>> population = initPopulation(testPool);
        double[] fit = evalAll(population);

        int bestIdx = argmax(fit);
        List<TestCase> bestEver = new ArrayList<>(population.get(bestIdx));
        double bestFitness = fit[bestIdx];

        for (int gen = 0; gen < generations; gen++) {
            List<List<TestCase>> next = new ArrayList<>(populationSize);
            double[] nextFit = new double[populationSize];

            // элитизм
            next.add(new ArrayList<>(bestEver));
            nextFit[0] = bestFitness;

            while (next.size() < populationSize) {
                List<TestCase> child1 = new ArrayList<>(tournament(population, fit));
                List<TestCase> child2 = new ArrayList<>(tournament(population, fit));

                boolean crossed = false;
                if (rng.nextDouble() < crossoverRate) {
                    oxCrossover(child1, child2);
                    crossed = true;
                }

                boolean mutated1 = false, mutated2 = false;
                if (rng.nextDouble() < mutationRate) { mutate(child1, mutationIntensity); mutated1 = true; }
                if (rng.nextDouble() < mutationRate) { mutate(child2, mutationIntensity); mutated2 = true; }

                int i = next.size();
                next.add(child1);
                nextFit[i] = (crossed || mutated1)
                        ? fitness(child1)
                        : fit[indexOf(population, child1)];

                if (next.size() < populationSize) {
                    int j = next.size();
                    next.add(child2);
                    nextFit[j] = (crossed || mutated2)
                            ? fitness(child2)
                            : fit[indexOf(population, child2)];
                }
            }

            population = next;
            fit = nextFit;

            int idx = argmax(fit);
            if (fit[idx] > bestFitness) {
                bestFitness = fit[idx];
                bestEver = new ArrayList<>(population.get(idx));
            }
        }

        return bestEver;
    }

    private List<List<TestCase>> initPopulation(List<TestCase> pool) {
        List<List<TestCase>> pop = new ArrayList<>(populationSize);
        for (int i = 0; i < populationSize; i++) {
            List<TestCase> ind = new ArrayList<>(pool);
            Collections.shuffle(ind, rng);
            pop.add(ind);
        }
        return pop;
    }

    private double[] evalAll(List<List<TestCase>> population) {
        double[] f = new double[population.size()];
        for (int i = 0; i < population.size(); i++)
            f[i] = fitness(population.get(i));
        return f;
    }

    private List<TestCase> tournament(List<List<TestCase>> population, double[] fit) {
        int best = rng.nextInt(population.size());
        for (int i = 1; i < TOURNAMENT_SIZE; i++) {
            int candidate = rng.nextInt(population.size());
            if (fit[candidate] > fit[best]) best = candidate;
        }
        return population.get(best);
    }

    /**
     * OX (Order Crossover): сохраняет валидную перестановку.
     * Сегмент [start, end] из p1 копируется в child1,
     * остаток заполняется из p2 в порядке появления. Аналогично для child2.
     */
    private void oxCrossover(List<TestCase> p1, List<TestCase> p2) {
        int size = p1.size();
        if (size < 2) return;

        int start = rng.nextInt(size);
        int end = rng.nextInt(size);
        if (start > end) { int tmp = start; start = end; end = tmp; }
        if (start == end) end = Math.min(end + 1, size - 1);

        List<TestCase> c1 = buildOxChild(p1, p2, start, end);
        List<TestCase> c2 = buildOxChild(p2, p1, start, end);
        p1.clear(); p1.addAll(c1);
        p2.clear(); p2.addAll(c2);
    }

    private List<TestCase> buildOxChild(List<TestCase> donor, List<TestCase> filler, int start, int end) {
        int size = donor.size();
        List<TestCase> child = new ArrayList<>(Collections.nCopies(size, null));
        Set<String> inSegment = new HashSet<>();

        for (int i = start; i <= end; i++) {
            child.set(i, donor.get(i));
            inSegment.add(donor.get(i).getId());
        }

        Queue<TestCase> rest = new ArrayDeque<>();
        for (TestCase t : filler)
            if (!inSegment.contains(t.getId())) rest.add(t);

        for (int i = 0; i < size; i++)
            if (child.get(i) == null) child.set(i, rest.poll());

        return child;
    }

    private void mutate(List<TestCase> ind, int intensity) {
        for (int k = 0; k < intensity; k++)
            Collections.swap(ind, rng.nextInt(ind.size()), rng.nextInt(ind.size()));
    }

    private int argmax(double[] arr) {
        int best = 0;
        for (int i = 1; i < arr.length; i++)
            if (arr[i] > arr[best]) best = i;
        return best;
    }

    private int indexOf(List<List<TestCase>> population, List<TestCase> ind) {
        for (int i = 0; i < population.size(); i++)
            if (population.get(i) == ind) return i;
        return 0;
    }

    @Override
    public String getName() { return "GeneticPrioritizer"; }
}
\end{lstlisting}

\begin{lstlisting}[language=Java, breaklines=true]
package com.tcp.algorithm;

import com.tcp.model.TestCase;
import java.util.List;

/*
интерфейс для всех алгоритмов приоритизации
позволяет единообразно вызывать любые стратегии
 */
public interface Prioritizer {
    /**
     * основной метод
     * принимает неупорядоченный пул тестов и возвращает список, отсортированный по приоритету
     *
     * @param testPool  - исходный список тестов
     * @return упорядоченный список тестов
     */
    List<TestCase> prioritize(List<TestCase> testPool);

    /**
     * возвращает название алгоритма для логирования и отчетов
     * @return имя метода приоритизации
     */
    String getName();

}
\end{lstlisting}

\begin{lstlisting}[language=Java, breaklines=true]
package com.tcp.algorithm;

import com.tcp.model.TestCase;

import java.util.*;

/**
 * Random Search -- случайная перестановка тестов.
 * Используется как baseline для сравнения с оптимизационными алгоритмами.
 *
 * Для воспроизводимости результатов используется фиксированный seed.
 * APFD усредняется по нескольким прогонам в Main.
 */
public class RandomPrioritizer implements Prioritizer {

    private final Random rng;

    // seed для воспроизводимости
    public RandomPrioritizer(long seed) {
        this.rng = new Random(seed);
    }

    @Override
    public List<TestCase> prioritize(List<TestCase> testPool) {
        List<TestCase> result = new ArrayList<>(testPool);
        Collections.shuffle(result, rng);
        return result;
    }

    @Override
    public String getName() { return "Random"; }
}

\end{lstlisting}

\begin{lstlisting}[language=Java, breaklines=true]
    package com.tcp.data;

import com.tcp.model.TestCase;

import java.io.*;
import java.util.*;
import java.util.stream.Collectors;

public class DataLoader {

    public static class FaultMatrixData {
        public final List<TestCase> tests;
        public final Map<String, Set<String>> testToFaults;
        public final int totalFaults;

        public FaultMatrixData(List<TestCase> tests, Map<String, Set<String>> testToFaults, int totalFaults) {
            this.tests = tests;
            this.testToFaults = testToFaults;
            this.totalFaults = totalFaults;
        }
    }

    // --- BigFaultMatrix: testId,0,1,0,... ---

    public static FaultMatrixData loadFaultMatrix(String filePath) throws Exception {
        List<TestCase> tests = new ArrayList<>();
        Map<String, Set<String>> testToFaults = new HashMap<>();
        int totalFaults = -1;

        for (String line : readLines(filePath)) {
            String[] parts = line.split(",");
            if (parts.length < 2) continue;
            if (totalFaults == -1) totalFaults = parts.length - 1;

            String testId = parts[0].trim();
            Set<String> killed = new HashSet<>();
            for (int i = 1; i < parts.length; i++)
                if ("1".equals(parts[i].trim())) killed.add("F" + (i - 1));

            double fp = totalFaults > 0 ? (double) killed.size() / totalFaults : 0.0;
            tests.add(new TestCase(testId, killed, fp, 100.0));
            testToFaults.put(testId, killed);
        }

        return new FaultMatrixData(tests, testToFaults, Math.max(totalFaults, 0));
    }

    public static FaultMatrixData createVariation(FaultMatrixData original, double flipRatio, long seed) {
        Random rng = new Random(seed);
        List<TestCase> newTests = new ArrayList<>();
        Map<String, Set<String>> newMap = new HashMap<>();

        for (TestCase t : original.tests) {
            Set<String> faults = new HashSet<>(original.testToFaults.get(t.getId()));
            for (int f = 0; f < original.totalFaults; f++) {
                if (rng.nextDouble() < flipRatio) {
                    String id = "F" + f;
                    if (!faults.add(id)) faults.remove(id);
                }
            }
            double fp = original.totalFaults > 0 ? (double) faults.size() / original.totalFaults : 0.0;
            newTests.add(new TestCase(t.getId(), faults, fp, 100.0));
            newMap.put(t.getId(), faults);
        }

        return new FaultMatrixData(newTests, newMap, original.totalFaults);
    }

    // --- Defects4J / tcp-methods ---
    // Struktura papki projectPath:
    //   coverage/matrix.csv          -- zagolovok: test,method1,...; stroki: testId,0,1,...
    //   mutants/kill.csv             -- MutantNo,[FAIL|TIME|EXC|LIVE|UNCOV]
    //   mutants/testMap.csv          -- TestNo,TestName,Runtime
    //   mutants/covMap.csv           -- TestNo,MutantNo
    //   cia/impact_probabilities.csv -- owner,method,descriptor,roc,fwd,impact

    public static FaultMatrixData loadDefects4JData(String projectPath) throws Exception {
        List<TestCase> tests = loadCoverage(projectPath + "/coverage/matrix.csv");
        Map<String, Set<String>> testToFaults = loadKills(
                projectPath + "/mutants/kill.csv",
                projectPath + "/mutants/testMap.csv",
                projectPath + "/mutants/covMap.csv",
                tests);
        tests = applyImpact(tests, projectPath + "/cia/impact_probabilities.csv");

        int totalFaults = (int) testToFaults.values().stream().flatMap(Set::stream).distinct().count();
        System.out.printf("  Tests: %d, mutants: %d%n", tests.size(), totalFaults);
        return new FaultMatrixData(tests, testToFaults, totalFaults);
    }

    private static List<TestCase> loadCoverage(String path) throws Exception {
        List<TestCase> tests = new ArrayList<>();
        List<String> methods = new ArrayList<>();

        try (BufferedReader br = new BufferedReader(new FileReader(new File(path)))) {
            String line;
            boolean header = true;
            while ((line = br.readLine()) != null) {
                line = line.trim();
                if (line.isEmpty()) continue;
                String[] parts = splitCsv(line);
                if (header) {
                    // первая строка - заголовок с именами методов
                    for (int i = 1; i < parts.length; i++)
                        methods.add(parts[i].replace("\"", "").trim());
                    header = false;
                    continue;
                }
                Set<String> covered = new HashSet<>();
                for (int i = 1; i < parts.length && i - 1 < methods.size(); i++)
                    if ("1".equals(parts[i].trim())) covered.add(methods.get(i - 1));
                tests.add(new TestCase(parts[0].trim(), covered, 0.0, 100.0));
            }
        }
        return tests;
    }

    private static Map<String, Set<String>> loadKills(String killPath, String testMapPath,
                                                      String covMapPath, List<TestCase> tests) throws Exception {
        Map<String, Set<String>> result = new HashMap<>();
        for (TestCase t : tests) result.put(t.getId(), new HashSet<>());

        // kill.csv: MutantNo -> status (FAIL/TIME/EXC = killed)
        Set<String> killed = new HashSet<>();
        for (String line : readLines(killPath)) {
            String[] p = line.split(",");
            if (p.length < 2) continue;
            String status = p[1].trim();
            if (status.equals("FAIL") || status.equals("TIME") || status.equals("EXC"))
                killed.add(p[0].trim());
        }

        // testMap.csv: TestNo -> TestName (className only, no method)
        Map<String, String> testNoToName = new HashMap<>();
        for (String line : readLines(testMapPath)) {
            String[] p = line.split(",");
            if (p.length < 2) continue;
            testNoToName.put(p[0].trim(), p[1].trim());
        }

        // coverage keys look like "org.jsoup.FormElementTest#methodName"
        // testMap has only "org.jsoup.FormElementTest" (class name, no method)
        // build index: className -> list of full testIds from coverage
        Map<String, List<String>> classToIds = new HashMap<>();
        for (String testId : result.keySet()) {
            String cls = testId.contains("#") ? testId.substring(0, testId.indexOf('#')) : testId;
            classToIds.computeIfAbsent(cls, k -> new ArrayList<>()).add(testId);
        }

        // covMap.csv: TestNo,MutantNo - assign killed mutants to tests
        for (String line : readLines(covMapPath)) {
            String[] p = line.split(",");
            if (p.length < 2) continue;
            String testNo = p[0].trim(), mutantNo = p[1].trim();
            if (!killed.contains(mutantNo)) continue;
            String className = testNoToName.get(testNo);
            if (className == null) continue;
            for (String testId : classToIds.getOrDefault(className, Collections.emptyList()))
                result.get(testId).add("M" + mutantNo);
        }

        return result;
    }

    private static List<TestCase> applyImpact(List<TestCase> tests, String path) throws Exception {
        if (!new File(path).exists()) return tests;

        // Нормализуем ключи CIA: org/jsoup/nodes/FormElement,<init>,... -> org.jsoup.nodes.FormElement#init
        Map<String, Double> scores = new HashMap<>();
        for (String line : readLines(path)) {
            String[] p = splitCsv(line);
            if (p.length < 6) continue;
            try {
                String className = p[0].trim().replace('/', '.');
                // убираем угловые скобки: <init> -> init, <clinit> -> clinit
                String methodName = p[1].trim().replace("<", "").replace(">", "");
                String key = className + "#" + methodName;
                double score = Double.parseDouble(p[5].trim());
                scores.merge(key, score, Math::max);
            } catch (NumberFormatException ignored) {}
        }

        // Нормализуем ключи coverage до того же формата:
        // "org.jsoup.parser$Parser#Parser(org.jsoup.parser.TreeBuilder):24" -> "org.jsoup.parser.Parser#Parser"
        return tests.stream().map(t -> {
            double sum = 0.0; int count = 0;
            for (String m : t.getCoveredMethods()) {
                String key = normalizeCoverageKey(m);
                if (scores.containsKey(key)) { sum += scores.get(key); count++; }
            }
            return new TestCase(t.getId(), t.getCoveredMethods(),
                    count > 0 ? sum / count : 0.0, t.getExecutionTimeMs());
        }).collect(Collectors.toList());
    }

    // "org.jsoup.parser$Parser#Parser(org.jsoup.parser.TreeBuilder):24" -> "org.jsoup.parser.Parser#Parser"
    private static String normalizeCoverageKey(String raw) {
        String s = raw.replaceAll(":\\d+$", "");   // убираем :24
        int paren = s.indexOf('(');
        if (paren >= 0) s = s.substring(0, paren); // убираем параметры
        return s.replace('$', '.');                 // $ -> . для вложенных классов
    }


    // --- Synthetic data ---

    /**
     * Генерирует синтетические данные с независимым риском и покрытием.
     *
     * Тесты делятся на три группы чтобы GA-FP+Div имел реальное преимущество:
     *   (30%): большое покрытие, мало фолтов -- жадный выбирает их первыми.
     *   (20%): малое покрытие, высокий риск и много ключевых фолтов.
     *   Обычные (50%): покрытие и фолты пропорциональны density.
     */
    public static FaultMatrixData generateSynthetic(int numTests, int numFaults, double density) {
        List<TestCase> tests = new ArrayList<>();
        Map<String, Set<String>> map = new HashMap<>();
        Random rng = new Random(42);

        // первые 40% фолтов считаются "рискованными"
        int riskFaults = (int)(numFaults * 0.4);

        for (int t = 0; t < numTests; t++) {
            Set<String> faults = new HashSet<>();
            double fp;
            int execTime = 50 + rng.nextInt(150);
            double roll = rng.nextDouble();

            if (roll < 0.30) {
                // "Ловушка": много обычных фолтов, но не рискованных -- низкий риск
                for (int f = riskFaults; f < numFaults; f++)
                    if (rng.nextDouble() < density * 2.5) faults.add("F" + f);
                fp = 0.05 + rng.nextDouble() * 0.05;

            } else if (roll < 0.50) {
                // "Жемчужина": мало покрытия, но попадает в рискованные фолты
                for (int f = 0; f < riskFaults; f++)
                    if (rng.nextDouble() < density * 1.5) faults.add("F" + f);
                fp = 0.6 + rng.nextDouble() * 0.4;

            } else {
                // Обычный тест
                for (int f = 0; f < numFaults; f++)
                    if (rng.nextDouble() < density) faults.add("F" + f);
                fp = numFaults > 0 ? (double) faults.size() / numFaults : 0.0;
            }

            TestCase tc = new TestCase("T" + t, faults, fp, execTime);
            tests.add(tc);
            map.put(tc.getId(), new HashSet<>(faults));
        }

        return new FaultMatrixData(tests, map, numFaults);
    }

    // --- Utils ---

    // Read file line by line, skip empty lines and header (first line)
    private static List<String> readLines(String path) throws Exception {
        List<String> lines = new ArrayList<>();
        File f = new File(path);
        if (!f.exists()) return lines;
        try (BufferedReader br = new BufferedReader(new FileReader(f))) {
            String line;
            boolean header = true;
            while ((line = br.readLine()) != null) {
                line = line.trim();
                if (line.isEmpty()) continue;
                if (header) { header = false; continue; }
                lines.add(line);
            }
        }
        return lines;
    }

    // Split CSV line respecting commas inside quotes
    private static String[] splitCsv(String line) {
        List<String> tokens = new ArrayList<>();
        StringBuilder sb = new StringBuilder();
        boolean inQuotes = false;
        for (char c : line.toCharArray()) {
            if (c == '"') inQuotes = !inQuotes;
            else if (c == ',' && !inQuotes) { tokens.add(sb.toString()); sb.setLength(0); }
            else sb.append(c);
        }
        tokens.add(sb.toString());
        return tokens.toArray(new String[0]);
    }
}
\end{lstlisting}

\begin{lstlisting}[language=Java, breaklines=true]
package com.tcp.fitness;

import com.tcp.model.TestCase;

import java.util.HashSet;
import java.util.List;
import java.util.Set;

/**
 * Фитнесс-функция для генетического алгоритма.
 *
 * Оценивает качество порядка тестов по трём компонентам, каждая в [0, 1]:
 *   coverage -- доля покрытых уникальных методов от всех в пуле
 *   faultProneness -- средний риск первых TOP_K тестов
 *   diversity -- среднее расстояние Жаккара между соседними парами в топ-TOP_K
 *
 * Итоговый фитнесс = взвешенная сумма трёх компонент (веса задаются в конструкторе).
 */
public class GaFitnessFunction {

    // число тестов в начале последовательности, учитываемых для risk и diversity
    private static final int TOP_K = 10;

    private final double weightCoverage;
    private final double weightFaultProneness;
    private final double weightDiversity;

    public GaFitnessFunction(double weightCoverage, double weightFaultProneness, double weightDiversity) {
        this.weightCoverage = weightCoverage;
        this.weightFaultProneness = weightFaultProneness;
        this.weightDiversity = weightDiversity;
    }

    // Вычисляет фитнесс последовательности.
    // totalMethods вычисляется каждый раз из пула - безопасно при смене проектов.
    public double calculateFitness(List<TestCase> seq) {
        if (seq.isEmpty()) return 0.0;

        Set<String> all = new HashSet<>();
        for (TestCase t : seq) all.addAll(t.getCoveredMethods());
        int totalMethods = Math.max(all.size(), 1);

        return weightCoverage      * coverage(seq, totalMethods)
                + weightFaultProneness * faultProneness(seq)
                + weightDiversity      * diversity(seq);
    }

    // Доля уникальных покрытых методов от totalMethods
    private double coverage(List<TestCase> seq, int totalMethods) {
        Set<String> covered = new HashSet<>();
        for (TestCase t : seq) covered.addAll(t.getCoveredMethods());
        return (double) covered.size() / totalMethods;
    }

    // Средняя fault-proneness первых TOP_K тестов
    private double faultProneness(List<TestCase> seq) {
        int k = Math.min(seq.size(), TOP_K);
        double sum = 0.0;
        for (int i = 0; i < k; i++) sum += seq.get(i).getFaultProneness();
        return sum / k;
    }

    // Среднее расстояние Жаккара между соседними парами в топ-TOP_K
    private double diversity(List<TestCase> seq) {
        int k = Math.min(seq.size(), TOP_K);
        if (k < 2) return 0.0;
        double sum = 0.0;
        for (int i = 0; i < k - 1; i++)
            sum += jaccard(seq.get(i).getCoveredMethods(), seq.get(i + 1).getCoveredMethods());
        return sum / (k - 1);
    }

    // Расстояние Жаккара: 0 = одинаковые, 1 = полностью разные множества
    private double jaccard(Set<String> a, Set<String> b) {
        if (a.isEmpty() && b.isEmpty()) return 0.0;
        if (a.isEmpty() || b.isEmpty()) return 1.0;

        // считаем пересечение без создания нового множества
        int intersection = 0;
        for (String s : a) if (b.contains(s)) intersection++;
        int union = a.size() + b.size() - intersection;

        return 1.0 - (double) intersection / union;
    }
}
\end{lstlisting}

\begin{lstlisting}[language=Java, breaklines=true]
package com.tcp.metric;

import com.tcp.model.TestCase;

import java.util.*;

/**
 * Метрика APFD (Average Percentage of Faults Detected).
 * Оценивает скорость обнаружения дефектов: чем ближе к 1, тем раньше находятся мутанты.
 *
 * Формула: APFD = 1 - sumTF / (n * m) + 1 / (2n)
 *   n -- число тестов
 *   m -- число мутантов
 *   sumTF -- сумма позиций первых тестов, обнаруживших каждый мутант
 */
public class ApfdCalculator {

    public static double calculate(List<TestCase> orderedTests,
                                   Map<String, Set<String>> testToMutants,
                                   int totalMutants) {
        int n = orderedTests.size();
        int m = totalMutants;
        if (n == 0 || m == 0) return 0.0;

        Set<String> detected = new HashSet<>();
        double sumTF = 0.0;

        for (int i = 0; i < n; i++) {
            for (String mutant : testToMutants.getOrDefault(orderedTests.get(i).getId(), Collections.emptySet())) {
                if (detected.add(mutant)) sumTF += (i + 1);
            }
        }

        return 1.0 - sumTF / ((double) n * m) + 1.0 / (2.0 * n);
    }
}
\end{lstlisting}

\begin{lstlisting}[language=Java, breaklines=true]
package com.tcp.metric;

import com.tcp.model.TestCase;

import java.util.*;

/**
 * Метрика APFDc (APFD with Cost) - APFD с учётом времени выполнения тестов.
 *
 * В отличие от APFD, взвешивает обнаружение мутантов по кумулятивному времени:
 * мутант, найденный быстрым тестом в начале, ценится выше чем найденный медленным в конце.
 *
 * Формула: APFDc = sum(totalTime - cumTime_i + execTime_i / 2) / (totalTime * m)
 *   totalTime -- суммарное время всех тестов
 *   cumTime_i -- кумулятивное время до конца теста i
 *   m -- число мутантов
 */
public class ApfdCostCalculator {

    public static double calculateWithCost(List<TestCase> orderedTests,
                                           Map<String, Set<String>> testToMutants,
                                           int totalMutants) {
        if (orderedTests.isEmpty() || totalMutants == 0) return 0.0;

        double totalTime = orderedTests.stream().mapToDouble(TestCase::getExecutionTimeMs).sum();
        Set<String> detected = new HashSet<>();
        double weightedSum = 0.0;
        double cumTime = 0.0;

        for (TestCase test : orderedTests) {
            cumTime += test.getExecutionTimeMs();
            for (String mutant : testToMutants.getOrDefault(test.getId(), Collections.emptySet())) {
                if (detected.add(mutant))
                    weightedSum += totalTime - cumTime + test.getExecutionTimeMs() / 2.0;
            }
        }

        return weightedSum / (totalTime * totalMutants);
    }
}
\end{lstlisting}

\begin{lstlisting}[language=Java, breaklines=true]
    package com.tcp;

import com.tcp.algorithm.AdditionalPrioritizer;
import com.tcp.algorithm.GeneticPrioritizer;
import com.tcp.algorithm.Prioritizer;
import com.tcp.algorithm.RandomPrioritizer;
import com.tcp.data.DataLoader;
import com.tcp.data.DataLoader.FaultMatrixData;
import com.tcp.fitness.GaFitnessFunction;
import com.tcp.metric.ApfdCalculator;
import com.tcp.metric.ApfdCostCalculator;
import com.tcp.model.TestCase;

import java.io.*;
import java.time.LocalDateTime;
import java.time.format.DateTimeFormatter;
import java.util.*;

/**
 * Точка входа бенчмарка приоритизации регрессионных тестов.
 *
 * Алгоритмы:
 *   Random      -- случайная перестановка (среднее по 5 прогонам)
 *   Additional  -- жадный baseline по покрытию
 *   GA-Base     -- ГА только с покрытием (coverage weight = 1.0)
 *   GA-FP+Div   -- ГА с fault-proneness и диверсификацией (0.4 / 0.4 / 0.2)
 *
 * Метрики: APFD, APFDc
 * Результаты сохраняются в results/results.csv и results/results.json
 */
public class Main {

    private static final String BFM_PATH    = "data/bigfaultmatrix/bigfaultmatrix.txt";
    private static final String RESULTS_DIR = "results";
    private static final String CSV_PATH    = RESULTS_DIR + "/results.csv";
    private static final String JSON_PATH   = RESULTS_DIR + "/results.json";
    private static final int    RANDOM_RUNS = 5;

    private static final List<String> DEFAULT_PROJECTS = List.of(
            "Jsoup/93", "Csv/16", "JacksonDatabind/112",
            "Time/27", "Codec/18", "JacksonCore/26", "Lang/65"
    );

    // Хранит все результаты текущей сессии для записи в файлы
    private static final List<ResultRow> allResults = new ArrayList<>();

    public static void main(String[] args) {
        initResultsDir();
        Scanner sc = new Scanner(System.in);
        System.out.println("=== TCP Benchmark ===");

        while (true) {
            System.out.println("\n1. BigFaultMatrix (4 вариации)");
            System.out.println("2. Defects4J");
            System.out.println("3. Синтетические данные");
            System.out.println("4. Все Defects4J проекты");
            System.out.println("0. Выход");
            System.out.print("Выбор: ");
            String choice = sc.nextLine().trim();

            switch (choice) {
                case "1": runBigFaultMatrix(); break;
                case "2":
                    System.out.print("Путь к проекту (например data/raw/Jsoup/93): ");
                    String path = sc.nextLine().trim();
                    runBenchmark(path, loadOrNull(() -> DataLoader.loadDefects4JData(path)));
                    break;
                case "3": runBenchmark("synthetic", runSynthetic(sc)); break;
                case "4": runAllProjects(sc); break;
                case "0":
                    saveResults();
                    sc.close();
                    return;
                default: System.out.println("Введите 0-4.");
            }

            System.out.print("\nНажмите Enter...");
            sc.nextLine();
        }
    }

    // --- Прогоны ---

    private static void runBigFaultMatrix() {
        FaultMatrixData original = loadOrNull(() -> DataLoader.loadFaultMatrix(BFM_PATH));
        if (original == null) return;

        System.out.println("\n=== BigFaultMatrix: 4 вариации ===");

        String[] labels = { "BFM-original", "BFM-flip10", "BFM-flip20", "BFM-flip30" };
        double[] flips  = { 0.0, 0.1, 0.2, 0.3 };
        long[]   seeds  = { 0L,  1L,  2L,  3L  };

        System.out.printf("%n%-22s | %-8s | %-8s | %-8s | %-8s | %s%n",
                "Вариация", "Random", "Additional", "GA-Base", "GA-FP+Div", "Лучший");
        System.out.println("-".repeat(82));

        for (int v = 0; v < 4; v++) {
            FaultMatrixData data = (v == 0)
                    ? original
                    : DataLoader.createVariation(original, flips[v], seeds[v]);

            int gens = adaptiveGens(data.tests.size());
            Map<String, double[]> metrics = collectMetrics(data, gens);

            double[] rnd = metrics.get("Random");
            double[] add = metrics.get("Additional");
            double[] gab = metrics.get("GA-Base");
            double[] gaf = metrics.get("GA-FP+Div");

            String best = bestName(new double[]{rnd[0], add[0], gab[0], gaf[0]},
                    new String[]{"Random", "Additional", "GA-Base", "GA-FP+Div"});

            System.out.printf("%-22s | %-8.4f | %-8.4f | %-8.4f | %-8.4f | %s%n",
                    labels[v], rnd[0], add[0], gab[0], gaf[0], best);

            allResults.add(new ResultRow(labels[v], data.tests.size(), data.totalFaults,
                    rnd[0], rnd[1], add[0], add[1], gab[0], gab[1], gaf[0], gaf[1]));
        }
        System.out.println("-".repeat(82));
    }

    private static FaultMatrixData runSynthetic(Scanner sc) {
        int tests   = readInt(sc,    "Тестов (100): ",   100);
        int faults  = readInt(sc,    "Фолтов (50): ",     50);
        double dens = readDouble(sc, "Плотность (0.1): ", 0.1);
        return DataLoader.generateSynthetic(tests, faults, dens);
    }

    private static void runBenchmark(String label, FaultMatrixData data) {
        if (data == null || data.tests.isEmpty() || data.totalFaults == 0) {
            System.err.println("Нет данных для бенчмарка.");
            return;
        }

        int gens = adaptiveGens(data.tests.size());
        System.out.printf("  Поколений ГА: %d%n", gens);

        Map<String, double[]> metrics = collectMetrics(data, gens);
        printResults(metrics);

        double[] rnd = metrics.get("Random");
        double[] add = metrics.get("Additional");
        double[] gab = metrics.get("GA-Base");
        double[] gaf = metrics.get("GA-FP+Div");

        allResults.add(new ResultRow(label, data.tests.size(), data.totalFaults,
                rnd[0], rnd[1], add[0], add[1], gab[0], gab[1], gaf[0], gaf[1]));
    }

    private static void runAllProjects(Scanner sc) {
        System.out.print("Базовый путь к данным (например data/raw): ");
        String base = sc.nextLine().trim();

        System.out.println("\n=== ПАКЕТНЫЙ ЗАПУСК ===");
        System.out.printf("%-25s | %-8s | %-8s | %-8s | %-8s | %s%n",
                "Проект", "Random", "Additional", "GA-Base", "GA-FP+Div", "Лучший");
        System.out.println("-".repeat(82));

        for (String project : DEFAULT_PROJECTS) {
            FaultMatrixData data = loadOrNull(() -> DataLoader.loadDefects4JData(base + "/" + project));

            if (data == null || data.tests.isEmpty() || data.totalFaults == 0) {
                System.out.printf("%-25s | ПРОПУЩЕН%n", project);
                continue;
            }

            int gens = adaptiveGens(data.tests.size());
            Map<String, double[]> metrics = collectMetrics(data, gens);

            double[] rnd = metrics.get("Random");
            double[] add = metrics.get("Additional");
            double[] gab = metrics.get("GA-Base");
            double[] gaf = metrics.get("GA-FP+Div");

            String best = bestName(new double[]{rnd[0], add[0], gab[0], gaf[0]},
                    new String[]{"Random", "Additional", "GA-Base", "GA-FP+Div"});

            System.out.printf("%-25s | %-8.4f | %-8.4f | %-8.4f | %-8.4f | %s%n",
                    project, rnd[0], add[0], gab[0], gaf[0], best);

            allResults.add(new ResultRow(project, data.tests.size(), data.totalFaults,
                    rnd[0], rnd[1], add[0], add[1], gab[0], gab[1], gaf[0], gaf[1]));
        }
        System.out.println("-".repeat(82));
    }

    // --- Сбор метрик ---

    // Запускает все алгоритмы и возвращает Map: algorithmName -> [apfd, apfdc]
    private static Map<String, double[]> collectMetrics(FaultMatrixData data, int gens) {
        Map<String, double[]> result = new LinkedHashMap<>();

        // Random: среднее по RANDOM_RUNS прогонам
        double sumApfd = 0, sumApfdc = 0;
        for (int r = 0; r < RANDOM_RUNS; r++) {
            List<TestCase> ordered = new RandomPrioritizer(r).prioritize(new ArrayList<>(data.tests));
            sumApfd  += ApfdCalculator.calculate(ordered, data.testToFaults, data.totalFaults);
            sumApfdc += ApfdCostCalculator.calculateWithCost(ordered, data.testToFaults, data.totalFaults);
        }
        result.put("Random", new double[]{ sumApfd / RANDOM_RUNS, sumApfdc / RANDOM_RUNS });

        // Additional, GA-Base, GA-FP+Div
        List<String> names = List.of("Additional", "GA-Base", "GA-FP+Div");
        List<Prioritizer> algorithms = new ArrayList<>();
        algorithms.add(new AdditionalPrioritizer());
        algorithms.add(new GeneticPrioritizer(100, gens, 0.8, 0.2, new GaFitnessFunction(1.0, 0.0, 0.0)));
        algorithms.add(new GeneticPrioritizer(100, gens, 0.8, 0.2, new GaFitnessFunction(0.4, 0.4, 0.2)));

        for (int i = 0; i < algorithms.size(); i++) {
            List<TestCase> ordered = algorithms.get(i).prioritize(new ArrayList<>(data.tests));
            double apfd  = ApfdCalculator.calculate(ordered, data.testToFaults, data.totalFaults);
            double apfdc = ApfdCostCalculator.calculateWithCost(ordered, data.testToFaults, data.totalFaults);
            result.put(names.get(i), new double[]{ apfd, apfdc });
            // сохраняем упорядоченный список для вывода топ-3
            orderedResults.put(names.get(i), ordered);
        }
        // Random топ-3: берём первый прогон
        orderedResults.put("Random",
                new RandomPrioritizer(0).prioritize(new ArrayList<>(data.tests)));

        return result;
    }

    // хранит упорядоченные списки тестов для вывода топ-3
    private static final Map<String, List<TestCase>> orderedResults = new LinkedHashMap<>();

    private static void printResults(Map<String, double[]> metrics) {
        System.out.println("\n--- РЕЗУЛЬТАТЫ ---");
        System.out.printf("%-20s | %-8s | %-8s%n", "Алгоритм", "APFD", "APFDc");
        System.out.println("-".repeat(42));

        String bestName = null; double bestApfd = -1;
        for (Map.Entry<String, double[]> e : metrics.entrySet()) {
            System.out.printf("%-20s | %-8.4f | %-8.4f%n", e.getKey(), e.getValue()[0], e.getValue()[1]);
            if (e.getValue()[0] > bestApfd) { bestApfd = e.getValue()[0]; bestName = e.getKey(); }
        }
        System.out.println("-".repeat(42));
        System.out.printf("Лучший: %s (APFD = %.4f)%n", bestName, bestApfd);

        // Топ-3 тестов для каждого алгоритма
        System.out.println("\n--- ТОП-3 ТЕСТОВ ---");
        for (Map.Entry<String, List<TestCase>> e : orderedResults.entrySet()) {
            System.out.printf("%n%s:%n", e.getKey());
            System.out.printf("  %-12s | %-8s | %-10s%n", "Тест", "Risk", "Coverage");
            System.out.println("  " + "-".repeat(36));
            List<TestCase> ordered = e.getValue();
            for (int i = 0; i < Math.min(3, ordered.size()); i++) {
                TestCase t = ordered.get(i);
                System.out.printf("  %-12s | %-8.3f | %d%n",
                        t.getId(), t.getFaultProneness(),
                        t.getCoveredMethods().size());
            }
        }
        orderedResults.clear();
    }

    // --- Сохранение результатов ---

    private static void saveResults() {
        if (allResults.isEmpty()) return;
        try {
            new File(RESULTS_DIR).mkdirs();
            saveCsv();
            saveJson();
            System.out.println("\nРезультаты сохранены:");
            System.out.println("  CSV:  " + CSV_PATH);
            System.out.println("  JSON: " + JSON_PATH);
        } catch (Exception e) {
            System.err.println("Ошибка сохранения: " + e.getMessage());
        }
    }

    private static void saveCsv() throws Exception {
        // Используем ; как разделитель - Excel в русской локали ожидает именно его
        // sep= в первой строке подсказывает Excel какой разделитель использован
        boolean exists = new File(CSV_PATH).exists();
        try (PrintWriter pw = new PrintWriter(
                new OutputStreamWriter(new FileOutputStream(CSV_PATH, true), "UTF-8"))) {
            if (!exists) {
                pw.println("sep=;");  // директива для Excel - автоматически выберет разделитель
                pw.println("timestamp;dataset;tests;faults;" +
                        "random_apfd;random_apfdc;" +
                        "additional_apfd;additional_apfdc;" +
                        "gabase_apfd;gabase_apfdc;" +
                        "gafpdiv_apfd;gafpdiv_apfdc");
            }
            String ts = LocalDateTime.now().format(DateTimeFormatter.ofPattern("yyyy-MM-dd HH:mm:ss"));
            for (ResultRow r : allResults)
                pw.printf("%s;%s;%d;%d;%.4f;%.4f;%.4f;%.4f;%.4f;%.4f;%.4f;%.4f%n",
                        ts, r.dataset, r.tests, r.faults,
                        r.randomApfd, r.randomApfdc,
                        r.addApfd, r.addApfdc,
                        r.gaBaseApfd, r.gaBaseApfdc,
                        r.gaFpApfd, r.gaFpApfdc);
        }
    }

    private static void saveJson() throws Exception {
        String ts = LocalDateTime.now().format(DateTimeFormatter.ofPattern("yyyy-MM-dd HH:mm:ss"));
        StringBuilder sb = new StringBuilder("[\n");
        for (int i = 0; i < allResults.size(); i++) {
            ResultRow r = allResults.get(i);
            sb.append("  {\n");
            sb.append(String.format("    \"timestamp\": \"%s\",\n", ts));
            sb.append(String.format("    \"dataset\": \"%s\",\n", r.dataset));
            sb.append(String.format("    \"tests\": %d,\n", r.tests));
            sb.append(String.format("    \"faults\": %d,\n", r.faults));
            sb.append(String.format("    \"Random\":     {\"apfd\": %.4f, \"apfdc\": %.4f},\n", r.randomApfd, r.randomApfdc));
            sb.append(String.format("    \"Additional\": {\"apfd\": %.4f, \"apfdc\": %.4f},\n", r.addApfd, r.addApfdc));
            sb.append(String.format("    \"GA-Base\":    {\"apfd\": %.4f, \"apfdc\": %.4f},\n", r.gaBaseApfd, r.gaBaseApfdc));
            sb.append(String.format("    \"GA-FP+Div\": {\"apfd\": %.4f, \"apfdc\": %.4f}\n", r.gaFpApfd, r.gaFpApfdc));
            sb.append(i < allResults.size() - 1 ? "  },\n" : "  }\n");
        }
        sb.append("]");
        try (PrintWriter pw = new PrintWriter(new FileWriter(JSON_PATH))) {
            pw.print(sb);
        }
    }

    // --- Вспомогательные методы ---

    private static int adaptiveGens(int n) {
        return Math.max(50, (int)(Math.log(n) / Math.log(2)) * 10);
    }

    private static String bestName(double[] apfds, String[] names) {
        int best = 0;
        for (int i = 1; i < apfds.length; i++) if (apfds[i] > apfds[best]) best = i;
        return names[best];
    }

    private static void initResultsDir() {
        new File(RESULTS_DIR).mkdirs();
    }

    private static FaultMatrixData loadOrNull(ThrowingSupplier<FaultMatrixData> supplier) {
        try { return supplier.get(); }
        catch (Exception e) { System.err.println("Ошибка загрузки: " + e.getMessage()); return null; }
    }

    private static int readInt(Scanner sc, String prompt, int defaultVal) {
        System.out.print(prompt);
        String s = sc.nextLine().trim();
        return s.isEmpty() ? defaultVal : Integer.parseInt(s);
    }

    private static double readDouble(Scanner sc, String prompt, double defaultVal) {
        System.out.print(prompt);
        String s = sc.nextLine().trim();
        return s.isEmpty() ? defaultVal : Double.parseDouble(s);
    }

    @FunctionalInterface
    private interface ThrowingSupplier<T> { T get() throws Exception; }

    // --- Модель строки результата ---

    private static class ResultRow {
        final String dataset;
        final int tests, faults;
        final double randomApfd, randomApfdc;
        final double addApfd, addApfdc;
        final double gaBaseApfd, gaBaseApfdc;
        final double gaFpApfd, gaFpApfdc;

        ResultRow(String dataset, int tests, int faults,
                  double randomApfd, double randomApfdc,
                  double addApfd, double addApfdc,
                  double gaBaseApfd, double gaBaseApfdc,
                  double gaFpApfd, double gaFpApfdc) {
            this.dataset = dataset; this.tests = tests; this.faults = faults;
            this.randomApfd = randomApfd; this.randomApfdc = randomApfdc;
            this.addApfd = addApfd; this.addApfdc = addApfdc;
            this.gaBaseApfd = gaBaseApfd; this.gaBaseApfdc = gaBaseApfdc;
            this.gaFpApfd = gaFpApfd; this.gaFpApfdc = gaFpApfdc;
        }
    }
}
\end{lstlisting}